\pdfoutput=1
\documentclass[11pt]{article}
\usepackage[margin=1in]{geometry}
\usepackage{amsmath,amssymb,amsthm}
\usepackage{enumitem}
\usepackage{hyperref}
\hypersetup{hidelinks,
  pdftitle={Split bad primes and the tuple structure of IUT: a fourth report},
  pdfauthor={Toshihisa Urahigashi}}

\theoremstyle{plain}
\newtheorem{proposition}{Proposition}
\newtheorem{corollary}[proposition]{Corollary}
\newtheorem{lemma}[proposition]{Lemma}
\theoremstyle{definition}
\newtheorem{question}{Question}
\newtheorem{remark}{Remark}
\newtheorem{program}{Programme}

\newcommand{\Fmod}{F_{\mathrm{mod}}}
\newcommand{\Vmod}{\mathbb{V}_{\mathrm{mod}}}
\newcommand{\Vbad}{\mathbb{V}^{\mathrm{bad}}}
\newcommand{\Vgood}{\mathbb{V}^{\mathrm{good}}}
\newcommand{\VV}{\mathbb{V}}
\newcommand{\Ind}[1]{\textup{(Ind#1)}}
\newcommand{\SA}{\textup{(SA)}}
\newcommand{\ord}{\operatorname{ord}}
\newcommand{\Norm}{\operatorname{N}}

\title{Split bad primes and the tuple structure of\\
inter-universal Teichm\"uller theory\\
\large A fourth report: what the source states, and what follows}
\author{Toshihisa Urahigashi\\
\small\texttt{123026.urahigashi.toshihisa@gmail.com}}
\date{Version 1.0.3 --- 8 September 2026\\[3pt]
\small ORCID \href{https://orcid.org/0009-0004-0460-6242}{\texttt{0009-0004-0460-6242}}\\[3pt]
\small DOI \href{https://doi.org/10.5281/zenodo.22648944}{10.5281/zenodo.22648944} (all versions)}

\begin{document}
\maketitle

\begin{abstract}
This report collects statements about inter-universal Teichm\"uller theory that
are either quoted from the source or proved outright, together with the
computations that accompany them.  It records: that a tensor packet admits
place tuples of mixed reduction type, and that selecting only nonsplit bad
primes does not remove them --- over $\mathbb Q(\sqrt5)$ a congruence family has
every odd bad prime split with exactly one bad place, and for a separate
explicit curve the nonsplit selection retains between $1.23\%$ and $1.24\%$ of
the selected Tate-divisor degree; that the general-stage parameter is
$\lambda=j^2\ord_p(\underline{\underline q})$ and that when every $e_{v_i}\le p-2$ ---
the condition under which Step (v) of \cite{IUT4} leaves a place out of the set $I^*$
it defines --- \cite[Prop.~1.4(iii)]{IUT4} yields the stage margin
$j^2s-(j+2)$; where the source places the evaluation, quoted verbatim from the
proof of \cite[Cor.~3.12]{IUT3}; two constraints on the indeterminacies, namely
that $\Ind2$ acts by a unit scalar and therefore leaves the holomorphic hull
unchanged, and that labels are not recoverable from a raw procession; and the
identity between the two forms in which \cite[Thm.~1.10]{IUT4} prints the
procession-normalized coefficient, which is $\ell(\ell+1)/12$ times the stage-one
coefficient --- a ratio that no statement below depends on, Proposition~\ref{p:chi}
being the statement about it.  Every assertion is quoted or proved,
the statement about $\Ind2$ is for a fixed \emph{nonarchimedean} tuple and the one
about labels for a fixed
stage, and no verdict on the correctness of the theory or on $abc$ is asserted.
\end{abstract}

\section{Scope and status of the statements}\label{s:scope}

This report accompanies the explicit initial-data construction in
\cite[Thm.~1]{First}, whose conditions are verified there; the expanded
verification is \cite[\S\S2--5]{Firstfull}.  That construction uses
\[
 F_0=\mathbb Q(\sqrt5),\quad \pi=16+3\sqrt5,\quad a=\pi^{29},\quad
 E_0:y^2=x(x-1)(x-a),\quad \ell=7,
\]
\[
 F=F_0(i,E_0[15]),\qquad K=F(E_0[7])=F_0(i,E_0[105]).
\]
All multiplicative places outside $2,7$ are selected as bad.  The global
cover, compatible section and cusp, and all selected local orders are
verified in the cited construction, with the warrants separated there: the geometric
constructions are proved and \emph{are not verified by any program}, while the
seventh-power-free clause that gives the selected orders at \emph{all} places is
certified by the exhaustive sieve of \cite[\S9]{First} together with the norm
identity.  The field of moduli is $F_0$; the
selected places above $211$ are one good place and one bad place, with
$e_{\mathrm g}=1$, $e_{\mathrm b}=105$ and
$\ord_{211}(\underline{\underline q}_{\mathrm b})=29/7$.

We use \textbf{[S]} for a cited source statement or a report of an author's
claim, \textbf{[P]} for a deduction with proof or a reproducible computation,
and \textbf{[H]} for an unproved identification or source reading.
Reporting a claim as [S] does not certify its truth.
\emph{Every assertion below is either quoted from a source or proved.}  We do not
claim more than that.  In particular we do not claim that no reading is involved
anywhere: writing $S(R)$ for ``the source asserts $R$'', $S(R)$ and $R$ are different
statements, and $S(R)\not\Rightarrow R$.  Where a consequence is drawn from a
quotation, what is proved is either a statement about the text or a mathematical
statement whose unproved inputs are cited theorems, named as \textbf{[S]} where they
are used --- Propositions \ref{p:ind2hull} and \ref{p:allsplit} are of the second kind.
Two scopes are stated
explicitly rather than left to the reader: Proposition \ref{p:ind2hull} is about a
\emph{fixed} nonarchimedean tuple, and the procession argument of \S\ref{s:indets} is
about a \emph{fixed} stage.  An earlier version guaranteed that nothing here was
conditional on any reading; that guarantee could not
stand as written (2026-09-08).
The source versions are specified in the bibliography.  In particular, a
reference to Step (v) means the proof of Theorem 1.10 of \cite{IUT4}.



\section{Place tuples and mixed reduction}\label{s:mech}

\begin{proposition}[\textup{[P]}]\label{p:dichotomy}
For a rational prime $p$ let $\VV_p=\{v\in\VV:v\mid p\}$ and let $A$ be a
label set with $|A|\ge2$.  The tuple set of the local tensor packet is
$\VV_p^A$.  Nonconstant place tuples exist if and only if $|\VV_p|\ge2$,
equivalently if $\Fmod$ has at least two places above $p$.
Tuples containing both chosen-good and chosen-bad places exist if and only
if both kinds of places occur in $\VV_p$.
\end{proposition}

\begin{proof}
By \cite[Definition 3.1(e)]{IUT1}, the section $\VV\to\Vmod$ restricts to a
bijection above each $p$.  Distributing the tensor product of direct sums in
\cite[Propositions 3.1, 3.2]{IUT3} gives
\[
 \bigotimes_{\alpha\in A}\bigoplus_{v\in\VV_p}V_{\alpha,v}
 =\bigoplus_{t\in\VV_p^A}\bigotimes_{\alpha\in A}V_{\alpha,t(\alpha)}.
\]
Each assertion now follows by choosing the values of a function $t:A\to\VV_p$.
No reduction hypothesis is needed for nonconstancy.
\end{proof}

Thus the mechanism involving mixed place tuples is absent when
$\Fmod=\mathbb Q$.  Over a quadratic field a nonsplit rational prime has a
single place, while a split prime has two; the split case can still have two
good places or two bad places.  Splitting alone does not supply mixed reduction.
We make no claim here about the full scope of the effective results in
\cite{MFHMP}.  Under Theorem 1.10 of \cite{IUT4}, chosen and actual reduction
types agree outside $2\ell$, as shown next.


\section{What changing the selected bad set can accomplish}\label{s:repair}

\begin{proposition}[the selection constraint; \textup{[P]}]\label{p:selection}
For a fixed initial datum satisfying the additional good-reduction hypothesis
of \cite[Theorem 1.10]{IUT4}, the chosen bad places outside $2\ell$ are exactly
the actual multiplicative places there.  Changing the chosen set alone cannot
eliminate a mixed profile at such a prime while retaining those hypotheses.
\end{proposition}

\begin{proof}
Definition 3.1(b) of \cite{IUT1} permits a nonempty subset of odd multiplicative
places; in particular it does not permit an actual-good place to be selected
as bad.  The extra hypothesis at the opening of \cite[Theorem 1.10, p.~22]{IUT4}
requires actual good reduction at each chosen-good place outside $2\ell$.
It therefore prohibits leaving an actual-bad place outside $2\ell$ unselected.
These two inclusions give the equality.  The distinction between the freedom
in the definition and the hypotheses for applying the theorem is essential.
\end{proof}

\begin{proposition}[a family with no odd nonsplit bad place; \textup{[P]}]
\label{p:allsplit}
Let $a=A-\sqrt5$, where $A\in\mathbb Z$, $A\equiv2\pmod5$ and
$A\equiv2\pmod{19}$.  For the Legendre curve over $F_0=\mathbb Q(\sqrt5)$,
every odd rational prime below a bad place splits in $F_0$, with exactly one
bad place above it.  The set of odd nonsplit bad places is empty.
\end{proposition}

\begin{proof}
The integral Legendre equation has
$\Delta=16a^2(a-1)^2$ and $c_4=16(a^2-a+1)$.  At an odd place it has good
reduction when $a\not\equiv0,1$, and multiplicative reduction when
$a\equiv0$ or $1$, since then $c_4$ is a unit.
For inert $p$, $\sqrt5\notin\mathbb F_p$, so $A-\sqrt5$ cannot be $0$ or $1$
in the residue field.  At the ramified prime $5$, its residue is $2$.
Now suppose that both places over a split odd $p$ were bad.  For a nonzero
root $s$ of $s^2=5$ in $\mathbb F_p$, the distinct values $A-s,A+s$ would
have to be $0,1$ in some order.  Hence $2A=1$ and $4s^2=1$ in
$\mathbb F_p$.  Thus $p\mid19$ and $A=10\pmod{19}$, contrary to the
hypothesis.  This proves the assertion.
\end{proof}

The separate Hensel/CRT existence family of \cite[Prop.~9.4]{Firstfull}
satisfies these congruences and has a bad place at $211$.  For its admissible
members, selecting only nonsplit odd bad places would select the empty set,
contrary even to the nonemptiness clause in \cite[Definition 3.1]{IUT1}.
Proposition~\ref{p:allsplit} is a statement about reduction for every stated
$A$; the additional initial-data conditions for the admissible subfamily
are not part of that proposition: \cite[Prop.~9.4]{Firstfull} supplies the arithmetic
--- the orders, the congruences and the seventh-power-free norm --- while
initial-$\Theta$-datum status for the family is claimed in the unnumbered paragraph
following it, on the strength of that draft's \S4, which we do not certify here.
We do not assert those
conditions for arbitrary powers $\pi^n$ or arbitrary $A$.

\subsection*{A different example: the numerical cost for $\pi^{29}$}

The following numbers concern $a=\pi^{29}=A_{29}+B_{29}\sqrt5$, not the
family $A-\sqrt5$.  Exact integer arithmetic gives
\[
 \gcd(A_{29},B_{29})=1,\qquad \gcd(A_{29}-1,B_{29})=3,
\]
\[
 R=\Norm(1-a)
 =25362449986390845063221249700759698691829298140087354726503101384580,
\]
and $v_2(R)=2$, $v_3(R)=2$, $v_5(R)=1$, $v_7(R)=0$.
At an inert odd $p$, the condition $a=0$ in the residue field requires
$p\mid A_{29},B_{29}$, which is impossible; the condition $a=1$ requires
$p\mid A_{29}-1,B_{29}$, leaving only $p=3$.  This prime is indeed inert and
bad.  At $5$, the residue of $a$ is $1$, giving the other nonsplit bad prime.
All other odd bad primes split.  Both places of a split prime cannot have
$a=0$ or both have $a=1$, by the gcds above.  If they had respectively
$a=0,1$, the prime would divide $\Norm(a)=211^{29}$, but the two residues
of $a$ at $211$ are $0,26$.  Thus these split bad primes have mixed reduction.

Define $H_{\rm all}$ as the degree-normalized Tate-divisor contribution from
bad places outside $2,7$, and $H_{\rm ns}$ as its nonsplit part, using
$1/[F_0:\mathbb Q]=1/2$.  The Tate orders at odd bad places are twice the
discrete orders of $a$ or $a-1$, so
\begin{equation}\label{e:heightcost}
 H_{\rm all}=29\log211+\log(R/4),\qquad H_{\rm ns}=2\log3+\log5=\log45.
\end{equation}
In the ramified case $5$, the residue degree is $1$ and the discrete order of
$a-1$ is $1$; the factor $1/2$ cancels the Tate order $2$, giving $\log5$.
The inert place at $3$ has residue degree $2$ and Tate order $2$, giving
$2\log3$ after normalization.  These are contributions to a selected divisor,
not full Weil heights.
The independent rational-series certificate in
\texttt{scripts/check\_round3.py} proves
\[
\begin{split}
309.021477380491&<H_{\rm all}<309.021477380492,\\
3.806662489770&<H_{\rm ns}<3.806662489771,\\
0.012318439876860&<H_{\rm ns}/H_{\rm all}<0.012318439876861.
\end{split}
\]
Discarding the split part would therefore discard more than $98.76\%$ of
this selected term, even if that operation were permitted.  In the
$A-\sqrt5$ family of Proposition~\ref{p:allsplit}, the corresponding
$H_{\rm ns}$ is instead zero.  The two statements must not be conflated.


\section{The logarithm lattice and the corrected stage margin}\label{s:hull}

\begin{lemma}[\textup{[P]}]\label{l:hull}
Let $p>2$ and $K_i/\mathbb Q_p$ be finite extensions with $e_i\le p-2$,
$0\le i\le j$.  In the tensor algebra put
$R_t=\bigotimes_{i,\mathbb Z_p}\mathcal O_{K_i}$,
$L_t=\bigotimes_{i,\mathbb Z_p}\log\mathcal O_{K_i}^\times$, and let
$\widetilde O_t$ be the product of integer rings of its field components.
For its ideal hull and degree-normalized logarithmic volume,
\[
 C_t(L_t)=\Bigl(\bigotimes_i\varpi_i\Bigr)\widetilde O_t,
 \qquad \mu^{\log}(C_t(L_t))=-\sum_i e_i^{-1}\log p,
\]
where $\varpi_i$ is a uniformizer.  If $L_t\subseteq C_t(U_t)$ then
$\mu^{\log}(C_t(U_t))\ge-\sum_i e_i^{-1}\log p$.
\end{lemma}

\begin{proof}
Under the stated bound, \cite[Proposition 1.2(i)]{IUT4} gives
$\log\mathcal O_{K_i}^\times=\mathfrak m_i=\varpi_i\mathcal O_{K_i}$.
Thus $L_t=(\bigotimes_i\varpi_i)R_t$ and its ideal hull has the displayed
form.  Each embedding of the tensor algebra into a field component assigns
valuation $\sum_i1/e_i$ to the generator, with $\ord_p(p)=1$.
The volume convention of \cite[Proposition 1.4(i)]{IUT4} yields the equality.
The inequality follows by monotonicity of the hull and its logarithmic volume.
No assertion that $L_t$ belongs to $U_t$ before taking the hull is required.
\end{proof}

Odd residue characteristic and tameness alone do not suffice.  For
$K=\mathbb Q_3(\zeta_3)$, $e=2=p-1$ is tame.  Put
$\varpi=\zeta_3-1$, $U_r=1+\mathfrak m^r$.
The quotient $U_1/U_2$ has order $3$ and is represented by $\mu_3$, so
$U_1=\mu_3U_2$.  The logarithm kills $\mu_3$, and logarithm and exponential
are inverse on $U_2$ and $\mathfrak m^2$ since their normalized valuations
are greater than $1/(p-1)$.  The remaining residue units are torsion.
Consequently $\log\mathcal O_K^\times=\mathfrak m^2\ne\mathfrak m$.
This also explains why exceeding $p-2$ should not be called the wild case:
that range contains tame extensions as well.

\begin{proposition}[the numerical margin; \textup{[P]}]\label{p:margin}
Under the hypotheses of Lemma~\ref{l:hull}, let the distinguished slot $j$
be bad and put $s=\ord_p(\underline{\underline q}_{t_j})$.
Step (v) of \cite[Theorem 1.10]{IUT4} uses $\lambda=j^2s$.
If the target hull contains $L_t$, its lower bound from the lemma exceeds
the expression $(-\lambda+d_I+1)\log p$ by
\[
 \bigl(j^2s-(j+2)\bigr)\log p.
\]
The difference is positive exactly when
\begin{equation}\label{e:threshold}
 j^2s>j+2,\qquad\text{equivalently}\qquad
 \ord_p(q_{t_j})>\frac{2\ell(j+2)}{j^2}.
\end{equation}
\end{proposition}

\begin{proof}
Here every extension is tame and $d_I=(j+1)-\sum_i1/e_i$.
The choice $I^*=\varnothing$ is admissible in
\cite[Proposition 1.4(iii)]{IUT4}.  Subtracting its numerical expression from
the lemma's bound gives
$-\sum_i1/e_i+j^2s-d_I-1=j^2s-(j+2)$.
The root is $q^{1/(2\ell)}$, proving the equivalent form.
\end{proof}

The proposition is a statement about the expression of
\cite[Prop.~1.4(iii)]{IUT4} at a given stage; whether that expression governs
the actual comparison region is a separate question, not decided here.  In the
explicit $\pi^{29}$ example,
$s=29/7$.  At $j=1$, target $(\mathrm g,\mathrm b)$ gives upper expression
$-226/105$, seed-hull volume $-106/105$ and margin $8/7$, after division by
$\log211$.  At $j=2$, source $(\mathrm b,\mathrm b,\mathrm g)$ can be
transported to target $(\mathrm b,\mathrm g,\mathrm b)$ by $(1\ 2)$,
fixing label $0$.  The corresponding values are
\[
\begin{aligned}
 &\lambda=116/7,\qquad d_I=208/105,\qquad a_I=107/105,\\
 &\text{upper}=-1427/105,\quad\text{seed}=-107/105,\quad
   \text{margin}=88/7.
\end{aligned}
\]
This checks a conditional use of an upper expression; it does not refute
Proposition 1.4 for the inputs covered by its hypotheses.

\subsection*{Which permutations matter}

For an admitted group $G\le\operatorname{Sym}(\{0,\ldots,j\})$, a target
tuple $t$ has a preimage with good distinguished slot precisely when
$t_k$ is good for some $k\in G\cdot j$.  Indeed, for
$(\sigma t)_i=t_{\sigma^{-1}i}$, one has
$(\sigma^{-1}t)_j=t_{\sigma j}$.
Thus protecting label $0$ does not block the $j=2$ example, but fixing $j$
does block this particular good-slot route when $t_j$ is bad.  A group may
fix $j$ and act nontrivially on other slots.  Our example therefore does not
require banning all nontrivial permutations as the only possible repair.
Nor does this orbit calculation by itself establish admission of any
permutation by the source or membership of its output in the actual union.


\section{Where the evaluation lives: two columns}\label{s:columns}

\textup{[S]} \cite[Cor.~3.12, proof, (iv)]{IUT3} states the aim of the passage:
\begin{quote}
``to apply the bi-coricity of the units [cf.\ Theorem 1.5, (iii), (iv)] in order to
compute the \emph{$0$-column} $\Theta$-pilot object in terms of the arithmetic
holomorphic structure of the \emph{$1$-column}.  In order to do this, one must work
with units that are vertically once-shifted [i.e., lie at $(n,m-1)$] relative to the
value group structures involved [i.e., which lie at $(n,m)$]''
\end{quote}
and resolves the conflict between ``vertical shift'' and ``impossibility of vertical
shift'' by working, on the $0$-column, with structures invariant under
$(0,m)\mapsto(0,m+1)$ --- which is what ``$(n,\circ)$'' denotes.  Subitem (iii)
explains why a single $\Theta^{\times\mu}_{\mathrm{LGP}}$-link is used: relating
units through one link and value groups through another would, the source says,
render it ``practically impossible to obtain conclusions that require one to relate
both \dots\ simultaneously'' --- that is subitem (iii).  \emph{The price of the
solution adopted is recorded separately, in subitem (iv)}: it ``may be implemented only
at the cost of admitting the `indeterminacy' constituted by the upper
semi-compatibility of $\Ind3$''.

\textup{[S]} Subitem (xi-b) places the output data:
\begin{quote}
``a construction of a collection of possibilities of output data contained in
$({}^{0,\circ}U^{\mathbb Q}\supseteq)\ {}^{0,\circ}U\xrightarrow{\ \sim\ }
{}^{1,\circ}U\ (\subseteq{}^{1,\circ}U^{\mathbb Q})$ --- where the isomorphism
`$\xrightarrow{\sim}$' arises from the permutation symmetries discussed in the final
portion of Theorem 3.11, (i) --- that satisfies the input prime-strip link (IPL) and
simultaneous holomorphic expressibility (SHE) properties''
\end{quote}
Subitem (xi-c) compares in the $1$-column, the collection being linked by (IPL), via
isomorphisms of $\mathcal F^\Vdash$-prime-strips, to the log-Kummer representation of
the $q$-pilot object at $(1,0)$; and (xi-d) says that the log-Kummer correspondence
must be applied ``in order to rectify the vertical shift/mismatch'' between the unit
portion of ${}^{1,0}\mathcal F^{\Vdash\times\mu}$ and the log-shells arising from it.
\emph{The iterates themselves are elsewhere}: \cite[Prop.~3.5(ii)(a)(2)]{IUT3} carries
the $m'$-th iterates for $m'\ge1$.  An earlier version of this paragraph attributed
``arbitrary iterates of the log-link'' to (xi-d); that phrase occurs nowhere in
(xi-b)--(xi-f), and joining the two places was a reading, not a quotation (internal check, 2026-09-08).

\textup{[P]} Three things follow from the quoted words themselves, without settling
what the passage establishes about the construction.
The comparison is not confined to one column: the possibilities lie in
${}^{0,\circ}U\xrightarrow{\sim}{}^{1,\circ}U$ and the comparison is made in the
$1$-column.  The units and the value groups sit at different vertical positions,
$(n,m-1)$ and $(n,m)$.  And the isomorphism identifying the two columns is
a \emph{poly}-isomorphism furnished by the permutation symmetries of
\cite[Thm.~3.11(i)]{IUT3} --- a collection of isomorphisms, so no single map is
selected here.

\begin{remark}[five things that may not be inferred]
\textup{[P]} A point lying in a log-shell is not thereby an output of the
$\Theta$-pilot.  The relevant statement, \cite[Prop.~3.5(ii)(a)]{IUT3}, is entitled
\emph{Upper Semi-compatibility} and asserts that the module in question ``contains
the images of the submodules of Galois invariants \dots\ of the \emph{groups of
units}'' --- an inclusion in one direction, about units, not an identification with
the $\Theta$-pilot output.  It also lists \emph{two} branches through which those
images arise: (1) the tensor product of the Kummer isomorphisms, and (2) their
pre-composite with the $m$-th iterates of the log-links \emph{for $m\ge1$}.  So the
zero-iterate branch is present, and one may not treat it as absent.

\textup{[S]} The source itself names the cost of the one-sided form.
\cite[Rem.~3.10.1(i)]{IUT3} says that the local monoids entering
$(^\dagger\mathcal F_{\mathrm{mod}})^\alpha$ ``are subject to a somewhat more
complicated `log-Kummer correspondence' \dots\ that revolves around `upper
semi-compatibility', i.e., in a word, \emph{one-sided inclusions, as opposed to
precise equalities}'', that ``the imprecise nature of such one-sided inclusions is
incompatible with the construction of'' that object, and that ``in particular, one
cannot construct log-link-compatible isomorphisms of Frobenioids'' for it --- in
contrast to $\mathcal F_{\mathrm{MOD}}$, for which the compatibility is precise.

\textup{[S]} \emph{The source states the counterpart in the same remark, and it must
be read with the above.}  \cite[Rem.~3.10.1(iii)]{IUT3} says that the local modules
of $\mathcal F_{\mathrm{mod}}$ --- ``i.e., local fractional ideals'' --- are, ``by
contrast'', ``well-suited to the computation of the étale-transport
indeterminacies''.  \emph{The same sentence begins by saying the opposite about the
other indeterminacy}: the additive nature of those modules renders $\Fmod$
ill-suited to the computation of Kummer-detachment indeterminacies.  It continues that
although ``various \emph{distortions} of these local
modules arise as a result of both \dots\ the local `upper semi-compatibility' of
Proposition 3.5, (ii), and \dots\ the \emph{discrepancy} between local holomorphic
and mono-analytic integral structures \dots, one may nevertheless compute --- i.e.,
if one takes into account the various distortions that occur, `estimate' --- the
global arithmetic degrees of objects of'' $\mathcal F_{\mathrm{mod}}$ ``by computing
log-volumes \dots, which are bi-coric''.  It then states: ``\emph{This computability
is precisely the topic of Corollary 3.12 below.}''  Consistently,
\cite[Thm.~3.11(ii)]{IUT3} records that the tensor-packet isomorphisms of its
subitem (a) --- not the Frobenioid isomorphisms of the preceding paragraph --- are
``[precisely!] compatible'' with the log-volumes, while the mutual compatibility of
the Frobenioid isomorphisms holds for the subscript ``MOD'' ``[but \emph{not}
between the isomorphisms that involve the subscript ``mod''! \dots]''.

\textup{[P]} So the one-sided inclusions above are \emph{acknowledged by the source}:
it locates them, names them as one of two sources of distortion, and asserts
computability at the level of log-volumes rather than at the level of Frobenioid
isomorphisms.  \emph{Whether that treatment is sufficient is a different question,
and not one decided here}; what is established is that pointing at the inclusions is
not by itself the exhibition of something the source has overlooked.  What the quoted passages do establish is where an
argument may and may not read an equality: at the level of the local modules the
relation is an inclusion, and a claim phrased as an identification of regions there
is not licensed by these statements.

\textup{[P]} Isomorphism class and degree do not determine a local representative.
In the category of \cite[Ex.~3.6(ii)]{IUT3} an object is a collection
$J=\{J_v\}_{v\in\mathbb V}$ of fractional ideals and, for $f\in F_{\mathrm{mod}}^\times$,
$f\cdot J=\{f\cdot J_v\}_v$ is again an object; the product formula makes the global
degree of $f\cdot J$ equal to that of $J$.  Taking $f=211^{-n}$ and $t$ a place above
$211$ gives $\bigcup_{n\ge0}211^{-n}J_t=K_t$.  So a fixed isomorphism class together
with a fixed global degree leaves the local representative unbounded.  The size of
the effect is exactly a height:

\begin{proposition}\label{p:reprheight}
\textup{[P]} Let $J=\{J_v\}$ be a fractional-ideal presentation of an arithmetic line
bundle over $F_0$ and $b\in F_0^\times$, so that $J$ and $bJ$ present the same bundle.
Write $E$ for the normalized global log-volume and $C(-)$ for the place-wise hull.
Then $E(bJ)=E(J)$, while
\[
 E\bigl(C(J\cup bJ)\bigr)-E(J)
 =\frac1{[F_0:\mathbb Q]}\sum_v n_v\log\max(1,|b|_v)=h(b),
\]
the absolute logarithmic Weil height of $b$.  For $b=211^{-n}$ over
$F_0=\mathbb Q(\sqrt5)$ this is $n\log211$.
\end{proposition}

\begin{proof}
At a finite place the two fractional ideals are comparable by inclusion and the hull
is the larger; at an infinite place they are balls with the same centre and the hull
is the larger.  So the local increment over $J_v$ is $\log\max(1,|b|_v)$, and the
weighted sum is $h(b)$.  The difference of the individual valuations is
$\sum_vn_v\log|b|_v/[F_0:\mathbb Q]=0$ by the product formula.  For $b=211^{-n}$ the
two places above $211$ each contribute $n\log211$ with weight $\tfrac12$ and the two
real places contribute $0$.
\end{proof}

\textup{[S]} \emph{The source fixes this by a definition, and says what happens
without it.}  \cite[Rem.~3.9.5(i)]{IUT3} defines the holomorphic hull \emph{in two
branches}, for a region already assumed to contain a relatively compact subset of
log-volume $>-\infty$: when the region ``is relatively compact'' the hull is the
smallest set of the form $\lambda\cdot\mathcal O$ containing it, with the components
of $\lambda$ nonzero; and otherwise --- ``\emph{If} $^\alpha U$ \dots\ \emph{is not
relatively compact, then we define the holomorphic hull of} $^\alpha U$ \dots\
\emph{to be} $I^{\mathbb Q}(^\alpha\mathcal F_{v_{\mathbb Q}})$'' --- it is the whole
ambient object.  \emph{The standing hypothesis is not decoration}: without it there
need be no smallest containing $\lambda\cdot\mathcal O$ at all, since
$U=\{0\}\subset\mathbb{Q}_p$ is compact and lies in every $p^{n}\mathbb{Z}_p$.
\textup{[P]} So relative compactness is not a convenience: in the second branch the
hull is still defined, but it gives no finite upper bound on the log-volume.  An
earlier version said the hull was defined ``only when'' the region is relatively
compact and that it then ``carries no information at all''; both overstated the
source.

\textup{[P]} \emph{The scalar value of a hull does not determine the value of the
transported hull.}  Work inside the good completion $K_{\mathrm g}$ of the datum of
\cite{Firstfull}.  \emph{That completion is unramified of degree $24$ over
$\mathbb{Q}_{211}$}, which we verify here rather than cite: $211$ splits in $F_0$, and
at the good place $\sqrt5\equiv146$, so $a=\pi^{29}\equiv26$; the reduced curve
$y^2=x(x-1)(x-26)$ over $\mathbb F_{211}$ has $a_{211}=24$ and $188$ points.  Since
$105$ is invertible there, N\'eron--Ogg--Shafarevich makes $K_{\mathrm g}/\mathbb Q_{211}$
unramified, of degree the order of Frobenius.  That order is the order of the companion
matrix because $X^2-24X+211$ is separable at each of $3,5,7$ --- its discriminant $-268$
is $2,2,5$ there --- so Frobenius is conjugate to it rather than scalar or unipotent.
The companion matrix of
$X^2-24X+211$ has order $4$ on $E[3]$, $3$ on $E[5]$ and $8$ on $E[7]$, hence $24$ on
$E[105]$, and $211\equiv3\pmod 4$ makes adjoining $i$ cost $2$, with
$\operatorname{lcm}(24,2)=24$.  \emph{The coincidence with $a_{211}=24$ is accidental.}
So
$e_{\mathrm g}=1<p-1$ and $\exp,\log$ are mutually inverse on $1+p\mathcal O$.
Normalize $\mu(\mathcal O)=1$ and write $\nu(U)=\tfrac1{24}\log\mu(U)$,
$E(U)=\nu(C(U))$ with $C$ the hull.

\begin{proposition}\label{p:nodescent}
Put $U_r=1+p^r\mathcal O$ for $r\ge1$.  Then $E(U_r)=0$ for every $r$, while
$E(\log U_r)=-r\log p$.  Consequently there is no function
$g:\mathbb{R}\to\mathbb{R}$ with $E(\log U)=g(E(U))$ for all these $U$ --- with or
without any linearity assumption.  Adding the raw volume does not repair this:
$U=1+p^3\mathcal O$ and $V=\exp(p)\,U$ satisfy
$\nu(U)=\nu(V)=-3\log p$ and $E(U)=E(V)=0$, yet $E(\log U)=-3\log p$ and
$E(\log V)=-\log p$.
\end{proposition}

\begin{proof}
$U_r$ contains $1$ and lies in $\mathcal O$, so its hull is $\mathcal O$ and
$E(U_r)=0$.  Since $r\ge1>e_{\mathrm g}/(p-1)$, $\log$ carries $U_r$ isometrically
onto $p^r\mathcal O$, which is already of the form $\lambda\mathcal O$; as
$\mu(p^r\mathcal O)=p^{-24r}$, $E(\log U_r)=-r\log p$.  The inputs are all $0$ and
the outputs are pairwise distinct, so no $g$ exists; and the outputs are unbounded
below, so no two-sided error term depending on the input alone repairs it either.
For the second pair, $\exp(p)$ is a unit, so $V\subseteq\mathcal O$ contains a unit
and $C(V)=\mathcal O$, while multiplication by a unit preserves $\mu$.  But
$\log V=p+p^3\mathcal O$ contains elements of valuation exactly $1$ and lies in
$p\mathcal O$, so $C(\log V)=p\mathcal O$.
\end{proof}

\textup{[S]} \emph{The source states the restriction that this respects.}
\cite[Rem.~3.9.6]{IUT3} says that the log-link compatibility ``amounts to a
coincidence of log-volumes --- \emph{not of arbitrary regions} that appear in the
domain and codomain of the log-link, but rather --- of certain types of
`\emph{sufficiently small}' regions''.  \textup{[P]} Proposition~\ref{p:nodescent}
is therefore not in tension with \cite{IUT3}: it exhibits what is lost if the
compatibility is read as a statement about hull scalars.  \emph{The point is not only
the smallness restriction}: the $U_r$ above are small, and $E$ is still constant on
them, because it is the passage to the hull --- not the size of the region --- that
discards what the raw log-volume $\nu$ records.

\textup{[S]} \emph{The pre-log-shell is the image of the invariants, not the
invariants of the quotient.}  \cite[Def.~1.1(i)]{IUT3} defines it as ``the image in
$\Psi^{\sim}_{{}^\dagger\mathcal F_v}$ of the submonoid of
$G_v({}^\dagger\Pi_v)$-invariants of $\Psi^{\times}_{{}^\dagger\mathcal F_v}$'', and
records that this ``constitutes a \emph{compact} topological module''; the log-shell
is that, times $(p_v^{*})^{-1}$.

\begin{proposition}\label{p:invimg}
\textup{[P]} Let $K_v$ be a finite extension of $\mathbb{Q}_p$, $\bar K_v$ an
algebraic closure, $G=G_{K_v}$, $\bar U=\mathcal O_{\bar K_v}^{\times}$, and
$\mu_\infty\subset\bar U$ the roots of unity.  Then
\[
 \log\,\operatorname{im}\bigl(\bar U^{G}\to\bar U/\mu_\infty\bigr)
   =\log\mathcal O_{K_v}^{\times},
 \qquad
 \log\bigl((\bar U/\mu_\infty)^{G}\bigr)=K_v .
\]
The first is a full lattice in $K_v$, hence compact; the second is not.
\end{proposition}

\begin{proof}
On units of an algebraic closure the kernel of $\log$ is exactly $\mu_\infty$: an
element lies in some finite extension, and after removing its finite-order residue
and passing to a high enough $p$-power it lies in a neighbourhood on which $\exp$ and
$\log$ are mutually inverse, so $\log z=0$ forces that power to be $1$.  It is
surjective onto $\bar K_v$: given $x$, choose $r$ with $p^rx$ in the domain of
convergence of $\exp$ inside a finite extension containing $x$, put $w=\exp(p^rx)$
and let $z$ be a $p^r$-th root of $w$ in $\bar K_v$; then $z$ is a unit and
$\log z=p^{-r}\log w=x$.  So $\log:\bar U/\mu_\infty\to\bar K_v$ is a
$G$-equivariant isomorphism, and taking invariants gives $K_v$ on the right.  The
image of the invariant units is $\log\mathcal O_{K_v}^{\times}$ by definition.
\end{proof}

\textup{[P]} So reading the pre-log-shell as the invariant part of the quotient
rather than as the image of the invariants replaces a compact lattice by the whole
local field, and any log-volume bound stated in terms of it becomes vacuous.  This is
not a defect of \cite{IUT3}, which takes the image and says so; it measures what the
other reading would cost.
\end{remark}

\section{Two constraints on the indeterminacies}\label{s:indets}

\textup{[P]} \emph{$\Ind2$ acts by a unit scalar} --- Proposition \ref{p:ind2hull}
below, at a fixed nonarchimedean tuple.  \emph{The classification of the isometries is
not ours}: \cite[Thm.~F, and Thm.~7.4]{MT} state that the natural
$\mathbb Z_p^{\times}\to\mathrm{Ism}(G_K)$ is bijective under the hypotheses recalled
next, and the authors summarize this as ``$\Ind2$ is relatively small''.  What is left
to this report is to check those hypotheses for a finite extension of $\mathbb Q_p$ and
to draw the consequence for the hull.  The identification is in the sources too:
\cite[Thm.~3.11(i)]{IUT3} describes $\Ind2$ as ``the indeterminacies induced by the
action of independent copies of $\mathrm{Ism}$ \dots\ on each of the direct summands of
the $j+1$ factors'', and the introduction of \cite{MT} states, \emph{before} Theorem F,
that ``$\Ind2$ consists of these isometries''.  \textup{[S]} \cite[Def.~7.3]{MT} defines
$\mathrm{Ism}(G_K)$ as the group of $G_K$-equivariant automorphisms of
$U_{K^{\mathrm{sep}}}^{\mu}$ that preserve the lattice $U_L^{\mu}$ for
\emph{each} finite separable $K\subseteq L\subseteq K^{\mathrm{sep}}$ --- not merely
one lattice --- and notes a natural injection
$\mathbb{Z}_p^{\times}\hookrightarrow\mathrm{Ism}(G_K)$.  Their Theorem 7.4 states
that \emph{if $G_k$ is not a pro-prime-to-$p$ group}, that injection is bijective;
their Remark 7.3.1 states that \emph{if $k$ is finite}, this $\mathrm{Ism}(G_K)$
is compatible with the one of \cite[Ex.~1.8(iv)]{IUT2}.  The authors themselves write that
``Theorem F implies that $\Ind2$ is relatively small''.

\textup{[P]} Both hypotheses hold in the case we use.  For $K$ a finite extension of
$\mathbb{Q}_p$ the residue field $k$ is finite, so Remark 7.3.1 applies; and
$G_k\cong\widehat{\mathbb{Z}}$ has $\mathbb{Z}_p$ as a quotient, so $G_k$ is not
pro-prime-to-$p$ and Theorem 7.4 applies.  Hence in this case
$\mathrm{Ism}(G_K)=\mathbb{Z}_p^{\times}$.

\begin{proposition}\label{p:ind2hull}
\textup{[P]} At a fixed \emph{nonarchimedean} tuple $t$, let $C_t$ denote the
holomorphic hull of \cite[Rem.~3.9.5(i)]{IUT3} in \emph{both} of its branches: the
smallest $\lambda\cdot\mathcal O$ containing the region when that region is relatively
compact, and the whole ambient object when it is not.  Assume that $U$ contains a
relatively compact subset whose log-volume is greater than $-\infty$, as required
by that definition.  Then for every $g$ in the group
above, $C_t(g_tU)=C_t(U)$ and $C_t\bigl(\bigcup_g g_tU\bigr)=C_t(U)$.
\end{proposition}

\begin{proof}
By the previous paragraph the action at $t$ is $g_t=u_t\,\mathrm{id}$ with
$u_t\in\mathbb{Z}_p^{\times}$, so each $g_t$ is an isometry: it preserves valuations
componentwise, hence preserves relative compactness and the family of sets
$\lambda\cdot\mathcal O$.

\emph{Relatively compact branch.}  If $U$ is relatively compact then so is $g_tU$, and
$U\subseteq\lambda\cdot\mathcal O$ if and only if $g_tU\subseteq\lambda\cdot\mathcal O$;
the two therefore have the same smallest such $\lambda\cdot\mathcal O$, so
$C_t(g_tU)=C_t(U)$.  The union $\bigcup_g g_tU$ lies in that same
$\lambda\cdot\mathcal O$ and contains $U$, so it is relatively compact with the same
hull.

\emph{Non-relatively-compact branch.}  If $U$ is not relatively compact then neither is
$g_tU$ nor the union, and all three hulls are the whole ambient object by definition.

\emph{An earlier version identified the hull with the $B_t$-module generated by the
region.}  That identification does not hold in general --- the module generated by a
region with full support is not of the form $\lambda\cdot\mathcal O$ --- and it is not
needed: the argument above uses only that $g_t$ is an isometry.
\end{proof}

\emph{In the relatively compact branch, this is proved using the same family of
$\lambda\mathcal O$ containers; in the other branch, it follows from the definition
of the hull.}  It concerns $\Ind2$ alone: it does not extend to
$\Ind1$, $\Ind3$, the input fibre or the $q$-comparison, and it does not say that an
arbitrary automorphism of $G_K$ acts as a ring isomorphism.

\textup{[P]} \emph{Labels are not recoverable from a raw procession.}  In a
procession whose entries are copies of one prime-strip, the internal symmetric group
of any stage acts while fixing every other stage, including the singleton stage: by
\cite[\S0]{IUT1} a capsule-full poly-morphism is a fixed injection together with
\emph{all} collections of isomorphisms, and by \cite[Def.~4.10]{IUT1} each arrow of a
procession is the collection of all such, so the automorphisms act stage-wise
independently, with the full symmetric group at each stage.  The only subsets
\emph{naturally selectable} --- that is, invariant under every automorphism of the raw
procession --- are therefore $\varnothing$ and the whole.  \textup{[S]} The source
states the conclusion directly: \cite[Prop.~6.9(i)]{IUT1} records that ``there are
precisely $n$ possibilities'' for the element of $|\mathbb F_\ell|$ to which a given
index of the $n$-capsule corresponds.  Hence one may not deduce, from the fixing of the
singleton $\mathcal D_0$, that the label $0$ is fixed inside the larger capsules; and
a coefficient computation for a subgroup does not by itself exhibit a repair, which
would require additional structure surviving in the actual weakened input together
with a proof that the choice is natural.  This concerns copies of one strip, not
processions with mutually non-isomorphic entries.

\section{The procession-normalized coefficient}\label{s:calib}

\textup{[S]} In Step (v) of the proof of \cite[Thm.~1.10]{IUT4} the weighted average
upper bound carries the term $-\frac{j^2}{2\ell}\log(\mathfrak q_{v_{\mathbb Q}})$ ---
the coefficient is \emph{negative}, and what is printed below are its magnitudes ---
at stage $j$ --- on the \emph{degree of the $q$-parameter divisor}, not on its
$2\ell$-th root: \cite[p.~23]{IUT4} records that $|\log(\underline{\underline q})|$
``is equal to'' $\tfrac1{2\ell}\log(\mathfrak q)$, so the coefficient on
$\log(\underline{\underline q})$ at stage $j$ is $j^2$.  An earlier version wrote
$\log(\underline{\underline q})$ in both places, conflating the two by a factor of
$2\ell$.  The source then says: ``it remains to compute the average over
$j\in\mathbb F_\ell^{*}$.  By averaging over $j\in\{1,\dots,\ell^{*}=\frac{\ell-1}2\}$
and applying (E1), (E2), we obtain the `procession-normalized upper bound'\,'' --- in
which the magnitude of the coefficient on $\log(\mathfrak q_{v_{\mathbb Q}})$ is printed first as
$\dfrac{(2\ell^{*}+1)(\ell^{*}+1)}{12\ell}$ and then as $\dfrac{\ell+1}{24}$.

\begin{proposition}\label{p:chi}
\textup{[P]} The two printed forms agree, and the average of the stage
coefficients equals $\dfrac{\ell(\ell+1)}{12}$ times the stage-one coefficient:
\[
 \frac1{\ell^{*}}\sum_{j=1}^{\ell^{*}}\frac{j^2}{2\ell}
 =\frac{(2\ell^{*}+1)(\ell^{*}+1)}{12\ell}
 =\frac{\ell+1}{24}
 =\frac{\ell(\ell+1)}{12}\cdot\frac1{2\ell}.
\]
Equivalently, the stage-one coefficient is
$\chi:=\dfrac{12}{\ell(\ell+1)}$ times the averaged one; at $\ell=7$,
$\chi=\tfrac3{14}$ and the averaged coefficient is $\tfrac13$.
\end{proposition}

\begin{proof}
$\sum_{j=1}^{m}j^2=m(m+1)(2m+1)/6$ with $m=\ell^{*}$ gives the first equality after
dividing by $\ell^{*}\cdot2\ell$.  Substituting $\ell^{*}=(\ell-1)/2$ gives
$2\ell^{*}+1=\ell$ and $\ell^{*}+1=(\ell+1)/2$, whence
$(2\ell^{*}+1)(\ell^{*}+1)/(12\ell)=(\ell+1)/24$.  The last equality is
$(\ell+1)/24=\bigl(\ell(\ell+1)/12\bigr)\cdot\bigl(1/(2\ell)\bigr)$.
\end{proof}

This is an identity between two quantities printed in the source; it is not a
statement about which region either coefficient is evaluated on.

\textup{[S]} In \cite[\S1.3, p.~4]{SS}, formula (1.5) is introduced as holding
``in essential approximation''.  Its coefficient on $\deg(q_E)$ is
$\frac{(\ell^{*}+1)(2\ell^{*}+1)}{12\ell}$, the form printed in
\cite[Thm.~1.10, Step (v)]{IUT4}.  The \emph{unnumbered} display immediately after
(1.5), obtained there by substitution into (1.4), is
\[
 \tfrac16\deg(q_E)\ \lessapprox\ \Bigl(1-\tfrac{12}{\ell(\ell+1)}\Bigr)^{-1}\cdot\mathfrak d(P).
\]
The following sentence in \cite[p.~4]{SS} retains an error, handled by adding a
term linear in $\ell$ on the right.  The display is not an error-free inequality.
\textup{[P]} \emph{That factor is $\chi$}: the ratio of the stage-one coefficient
$1/(2\ell)$ to the averaged $(\ell+1)/24$ is $12/(\ell(\ell+1))$.  Proposition
\ref{p:chi} thus identifies the averaged Tate-divisor coefficient under these
conventions.  It does not identify the full estimates, their positive or error
terms, or the regions to which they apply.

\textup{[S]} \emph{The source's own corresponding constant is a different one.}  In
assembling the local bounds of Steps (v)--(vii), \cite[Thm.~1.10, proof Step (viii),
p.~30]{IUT4} uses
$\frac{\ell+1}{4}\cdot\frac16\cdot\frac{12}{\ell^{2}}\ge\frac1{2\ell}$, that is
$12/\ell^{2}$, not $12/(\ell(\ell+1))$.  \textup{[P]} The two are not equal, and
$12/\ell^{2}>12/(\ell(\ell+1))$; the inequality just quoted has slack
$\frac{\ell+1}{2\ell^{2}}-\frac1{2\ell}=\frac1{2\ell^{2}}$.  This is an exact
comparison of these two scalar constants; the display in \cite{SS} still has the
approximation and error just stated.  \emph{What this records is that the averaging
coefficient is common to both and the final constant is not.  It decides nothing about
\cite[Cor.~3.12]{IUT3} or about the objection of} \cite{SS}.

Neither coefficient is used in any other statement of this report.

\begin{remark}[on explicit formulations]
\textup{[S]} \cite[\S4.10, eq.~(4.4)]{DH} writes the multiradial representation of
the $\Theta$-pilot region as
$U_\Theta:=\Ind2\bigl(\Ind1\bigl((\mathcal O_L(-P_\Theta))^{\Ind3}\bigr)\bigr)\subset L$,
and \cite[Rem.~4.10.1]{DH} states that ``the order of operations and formulas for the
indeterminacies do not appear in'' Tan's note or in \cite{IUT3}, \cite{IUT4}.  The
$\Ind3$ ingredient is moreover an assumed bound: \cite[\S4.9]{DH} writes ``we assume
a bound on $\mathcal O_L(-P_\Theta)^{\Ind3}$'', and the introduction to \cite[\S4]{DH}
describes it as ``a formula for an upper bound on the $\Ind3$ indeterminacies''.

\textup{[S]} \cite[p.~3]{DH} states, on the other hand, that the ordering in their
own (1.2) is deliberate, and reports that Mochizuki considers $\Ind1$ and $\Ind2$
together.  \textup{[P]} One consequence, and one non-consequence.  Exhibiting a map
inside an over-large group is not a lower bound for the actual indeterminacy, since
the bound is an upper one.  But a counterexample to a formulation of this shape is
not \emph{thereby} a counterexample to \cite{IUT3} either: that would require the
equivalence of the two formulations, which is not established here and which the two
passages above do not settle in either direction.  An earlier version inferred from
the first passage that the source is \emph{not committed} to the ordering; that does
not follow from a statement about where formulas appear .
\end{remark}

\section*{Corrections and reproducibility}

Earlier drafts omitted $j^2$ at stage $2$, producing the incorrect claimed
margin $1/7$; the corrected margin is $88/7$.  They also confused nonconstant
place tuples with mixed reduction, inferred independent probabilities
$2^{-r}$ from a splitting density, asserted the logarithm-lattice identity
for every tame odd place, and claimed that every nontrivial permutation
would have to be forbidden.  Those inferences are withdrawn; the proofs
and restrictions above replace them.  The earlier comparison of $1/3$
with $2/3$ also used an incorrect weighting.  The normalized bound
coefficients are not used in any statement above except
Proposition~\ref{p:chi}, which is the identity between them.

Corrections in the accompanying initial-data account include
$46080=2^{10}3^25$, the exceptional invariant $2^{14}31^3/5^3$,
selection of all multiplicative places outside $2,7$ including the place
over $59$, and the split-multiplicative argument after adjoining $i$.
Full rational $2$-torsion does not itself exclude a nonsplit twist.
The mistaken suggestion that all four exceptional invariants are algebraic
integers is also retained as an error record in the first report.

Every passage quoted above was checked against the English original of the version
listed in the bibliography.  \emph{The two steps are different}: the shipped
\path{scripts/note3_quotes.py} \emph{locates} a candidate passage by a word-sequence
match that drops mathematics and ignores punctuation, and reports twelve of the forty-one
entries as too short to locate at all; verbatim identity was then checked by reading.
The mechanical step used string match on the extracted text and, where text
extraction interleaves floats --- which it does on p.~150 of \cite{IUT3} --- by
reading the rendered page.  Two conventions of the extraction are worth recording,
since each produced a false mismatch before being accounted for: a page's text may
contain control characters inside a word, and the source uses double quotation marks
where a quotation reproduced here uses single ones.

The exact scripts released with the present report --- revised copies of those
published with \cite{First}, certifying the same numbers and additionally refusing to
run with assertions disabled --- verify the displayed
integer and rational arithmetic, including rigorous logarithm intervals and
the normalization identities.  They do not verify source interpretations.
The finite Lean model checks a distinction between invariance, orbit
membership and hulls of unions; it is not a model of the entire source theory.
Large language models were used in drafting, computations and source review.
No independent human verification or absence of further errors is asserted.


\begin{thebibliography}{First-full}
\bibitem[First]{First} T.~Urahigashi, \emph{An explicit initial $\Theta$-datum with
 mixed reduction at a split prime}, Zenodo, DOI 10.5281/zenodo.22537342, 2026.
\bibitem[SS]{SS} P.~Scholze, J.~Stix, \emph{Why $abc$ is still a conjecture},
 manuscript, July 16, 2018.
\bibitem[First-full]{Firstfull} T.~Urahigashi,
\emph{An explicit initial $\Theta$-datum with mixed reduction over a split prime,
and the local estimates of inter-universal Teichm\"uller theory IV},
version 0.1.3-draft (2026), accompanying manuscript
\texttt{paper.tex} in the accompanying release.
\bibitem[IUT1]{IUT1} S.~Mochizuki,
\emph{Inter-universal Teichm\"uller theory I}, RIMS Preprint 1756 (2012),
consulted author version, May 2020.
Published in Publ. Res. Inst. Math. Sci. \textbf{57} (2021), 3--207;
\href{https://doi.org/10.4171/PRIMS/57-1-1}{doi:10.4171/PRIMS/57-1-1}.
\bibitem[IUT3]{IUT3} S.~Mochizuki,
\emph{Inter-universal Teichm\"uller theory III}, RIMS Preprint 1758 (2012),
consulted author version, May 2020.
Published in Publ. Res. Inst. Math. Sci. \textbf{57} (2021), 403--626;
\href{https://doi.org/10.4171/PRIMS/57-1-3}{doi:10.4171/PRIMS/57-1-3}.
\bibitem[IUT4]{IUT4} S.~Mochizuki,
\emph{Inter-universal Teichm\"uller theory IV}, RIMS Preprint 1759 (2012),
consulted author version, April 2020.
Published in Publ. Res. Inst. Math. Sci. \textbf{57} (2021), 627--723;
\href{https://doi.org/10.4171/PRIMS/57-1-4}{doi:10.4171/PRIMS/57-1-4}.
\bibitem[DH]{DH} T.~Dupuy, A.~Hilado, \emph{The Statement of Mochizuki's Corollary
 3.12, Initial Theta Data, and the First Two Indeterminacies}, arXiv:2004.13228v1
 [math.NT], 28 April 2020, 35 pp.  \emph{Version~1 is cited deliberately.}  Under the
 same identifier, version~2 (19 June 2025, 20 pp.) is a different paper, on
 constructing initial theta data for elliptic curves over $\mathbb{Q}$; none of the
 four passages quoted above survives into it.
\bibitem[IUT2]{IUT2} S.~Mochizuki, \emph{Inter-universal Teichm\"uller theory II},
 Publ. Res. Inst. Math. Sci. \textbf{57} (2021), 209--401.
\bibitem[MT]{MT} A.~Minamide, S.~Tsujimura, \emph{Anabelian Geometry for Henselian
 Discrete Valuation Fields with Quasi-finite Residues}, RIMS-1973, Research Institute
 for Mathematical Sciences, Kyoto University, June 2023, 69 pp.
\bibitem[MFHMP]{MFHMP} S.~Mochizuki, I.~Fesenko, Y.~Hoshi, A.~Minamide,
W.~Porowski, \emph{Explicit estimates in inter-universal Teichm\"uller theory},
Kodai Math. J. \textbf{45} (2022), 175--236;
\href{https://doi.org/10.2996/kmj45201}{doi:10.2996/kmj45201}.
Consulted author versions: December 2021 and May 2022.
\end{thebibliography}
\end{document}
